\documentclass[10pt,a4paper]{article}
\usepackage[margin=2cm]{geometry}
\usepackage[T1]{fontenc}
\usepackage{lmodern}
\usepackage{booktabs,tabularx,array}
\usepackage[hidelinks]{hyperref}
\usepackage{microtype}
\setlength{\parindent}{0pt}
\setlength{\parskip}{5pt}
\title{UMT cluster inventory for MLIP nucleation research}
\author{MLIP-nucleation project}
\date{17 October 2026}
\begin{document}
\maketitle

\section{Scope and evidence}
Read-only SSH inspection was performed as \texttt{u3129} on
\texttt{umtnew.imm.uran.ru} (connection alias \texttt{umt.imm.uran.ru}).
The login node runs AlmaLinux 8.10 and Slurm 23.02.7.
No compute jobs were submitted, no VASP calculation or ML training was run,
and no installed environment was changed. Package versions below are installation
metadata, not proof of successful imports or GPU execution.

The versioned JSON observations record commands, timestamps in UTC, stdout,
stderr and return codes. They are stored under
\path{configs/observations/} with these suffixes after \texttt{2026-10-17-}:
\texttt{login-inventory.json}, \texttt{software-details.json},
\texttt{vasp-layout.json}, \texttt{user-software.json} and
\texttt{environment-activation.json}.
The collector is \path{scripts/inspect_cluster.py}.
Individual shell diagnostics must be checked as well as the final return code:
an earlier command in a compound probe can fail even if its final command succeeds.

\section{Slurm resources and access}
\begin{tabularx}{\linewidth}{@{}lrrXl@{}}
\toprule
Partition & Nodes & Slurm CPUs & Configured resources & Max time \\
\midrule
apollo & 36 & 1296 & 36 CPUs/node; memory 257000 or 385600 MB/node & 56 h \\
tesla & 33 & 442 & Heterogeneous CPUs; 15 K40m and 62 M2090 GPUs & 56 h \\
debug & 4 & 64 & Three CPU-only nodes and one node with 3 K40m GPUs & 90 min \\
v100 & 1 & 48 & \texttt{tesla-a101}: 2 A100, memory 254000 MB & 56 h \\
all & 33 & 442 & Same node set as \texttt{tesla}; default partition & 56 h \\
nompi & 69 & 1738 & Union of \texttt{apollo} and \texttt{tesla} & 56 h \\
\bottomrule
\end{tabularx}

Partition totals overlap and must not be added. Across unique reported nodes,
Slurm lists 74 nodes and 1850 CPUs. These are scheduler CPU counts, not a count of
physical cores. The A100 node has 2 threads/core, 24 physical cores and 48 Slurm CPUs.
The \texttt{debug} partition contains another 3 K40m GPUs, giving 18 K40m devices
across the unique nodes observed, plus 2 A100 and 62 M2090 devices configured in Slurm.
Configured devices on down nodes are not usable capacity.

At the snapshot, 11 \texttt{apollo} nodes, 8 \texttt{tesla} nodes and 3
\texttt{debug} nodes were idle. One node in each of \texttt{tesla} and
\texttt{debug} was down/unresponsive. The A100 node was mixed with 16 CPUs,
96 GiB of memory and one full A100 allocated; another A100 was not reported
allocated. It also advertises 12 GPU shards, which are sharing resources and
must not be interpreted as 12 additional GPUs. Device memory, driver version,
and practical whole-GPU versus shard scheduling remain unverified.
Availability is a timestamped observation, not a reservation or a guarantee.

The user belongs to groups \texttt{umt}, \texttt{umt9}, and \texttt{umtv100},
matching the displayed partition group restrictions. The default account is
\texttt{imet-2026-2}, with QoS \texttt{normal}. Its displayed per-user CPU limit
is 1024 and maximum submitted jobs per user is 2048. Older project associations
use QoS \texttt{restricted}, whose group CPU limit is zero.
The association wall-time limit is 365 days, but the shorter partition maximum
still applies. Displayed associations do not establish a remaining project
allocation balance; no spend budget has been chosen for this study.

Default memory is 1000 MB per allocated CPU and default wall time is 30 minutes.
Memory should be requested explicitly for training. Slurm uses
\texttt{select/cons\_tres}, \texttt{CR\_CORE\_MEMORY}, \texttt{task/cgroup},
and accounting enforcement \texttt{associations,limits,qos,safe}.
Partition billing weights make GPU requests materially more expensive than a CPU
request; account membership alone does not authorize an arbitrary allocation.

\section{VASP}
VASP was not found in the current PATH or the available module list.
\path{/opt/vasp} exists but is owned by \texttt{root:root} with mode 750;
listing and searching it are denied to this user. This does not prove that no
other installation exists. No usable VASP executable, version, MPI linkage or
licensed pseudopotential location has yet been established.

The next input needed is the path to an existing working VASP executable or
the user's established Slurm submission script. An administrator may need to
confirm access to the licensed installation. No POTCAR content was read or copied.

\section{MLIP environments and supporting software}
All environments below are under \path{/home/u3129/miniconda3/envs/}, resolving
through the home symlink to \path{/misc/home4/u3129/miniconda3/envs/}.

\begin{tabularx}{\linewidth}{@{}lX@{}}
\toprule
Environment & Observed package metadata \\
\midrule
\texttt{deepmd\_320\_a100} & Python 3.12.14; DeePMD-kit 3.2.0; PyTorch 2.12.0
(CUDA 13.0 build); TensorFlow 2.21.0; Open MPI 5.0.10; dpdata 1.1.0; e3nn 0.6.0.
Some CUDA toolkit components are 13.3; their presence does not establish driver compatibility. \\
\texttt{deepmd\_229\_cpu} & Python 3.10.13; DeePMD-kit 2.2.9; TensorFlow 2.9.0
(CPU); MPICH 4.3.2; LAMMPS 20230802.1. \\
\texttt{deepmd\_2210\_k40} & Python 3.10.13; DeePMD-kit 2.2.10; TensorFlow 2.9.0;
CUDA toolkit 10.2.89; cuDNN 7.6.5; MPICH 4.3.2. K40 runtime remains to be tested. \\
\texttt{lammps\_deepmd\_dpa4} & Python 3.12.14; DeePMD-kit 3.2.0; PyTorch 2.11.0;
ASE 3.29.0; LAMMPS 2025.7.22.2.0; MPICH 5.0.1.post1; CUDA 13 packages. \\
\texttt{lammps\_deepmd\_tf\_dpa4} & Python 3.12.14; DeePMD-kit 3.2.0;
TensorFlow 2.21.0; ASE 3.29.0; LAMMPS 2025.7.22.2.0; MPICH 5.0.2; CUDA 12 packages. \\
\bottomrule
\end{tabularx}

Activation, \texttt{python --version}, and resolution of \texttt{dp} and
\texttt{lmp} were verified for the first three environments. The last two were
inspected through package metadata only. MACE and NequIP distributions were not
found in the checked environments; installations elsewhere are not excluded.
Names containing \texttt{dpa4} are environment labels and do not verify model
architecture support, checkpoint availability, or successful model loading.

The default system Python is 3.6.8; the user's base Miniconda metadata reports
Python 3.14.7. Explicit activation or absolute environment paths are therefore
necessary. CUDA modules include 10.2, 11.8, 12.4, 12.6, 12.8 and 13.0, with
Open MPI 4.1.1 and Intel MPI/oneAPI 2025 modules also available. Module presence
does not certify compatibility with an environment's bundled MPI or CUDA runtime.
An older LAMMPS executable is on PATH at
\path{/home/u3129/Keller/SOFT/lammps-23Jun2022/src/lmp_mpi}.

\section{Storage and next checks}
The home directory resolves to \path{/misc/home4/u3129}; its filesystem had
about 24 TB free out of 40 TB at inspection. This is shared filesystem space,
not a personal quota. \path{~/_scratch} points to \path{/home1/u3129} and
\path{~/_scratch2} to \path{/home6/u3129}. The displayed quota for
\texttt{stg06:/home6} was 100 GB with about 15.7 GB used; the quota for the
current home filesystem was not established. Retention and backup policies
for these paths remain unknown. A production work root has not been selected.

Next steps are to locate accessible VASP and inspect its linkage, then run a
bounded Slurm environment check with an explicitly agreed resource budget.
For A100 this should record \texttt{nvidia-smi}, device memory and driver,
library imports, CUDA device discovery, and a tiny GPU operation.
K40 requires a separate check using its existing legacy environment.
A scientific VASP validation additionally needs a licensed potential and
a documented small input system. These runtime checks are not claimed complete.

\end{document}
